% !TEX root = ../main.tex

\chapter{Introduction}
\label{ch:introduction}

Start each chapter in its own file and include it from \texttt{main.tex}. This
keeps the root file focused on document structure and package configuration.

\section{Cross References}
\label{sec:cross-references}

Use labels for chapters, sections, figures, tables, equations, and appendix
material. For example, chapter~\ref{ch:examples} shows the figure and table
patterns, while appendix~\ref{app:additional-material} contains supporting
material. Citations are handled by \texttt{biblatex}; here is a textual citation
with \textcite{knuth1984texbook} and a parenthetical citation \cite{lamport1994latex}.

\begin{definition}[Reusable command]
A reusable command is a short macro in \texttt{macros.tex} that captures
notation used throughout the thesis.
\end{definition}

For example, \(\E[x \sim p]{x^2}\) is defined once in \texttt{macros.tex} and
can then be used consistently in every chapter.
